\makeatletter
\@ifundefined{keepfig}{\newenvironment{keepfig}{\par\edef\keepfig@pi{\the\parindent}\noindent
  \begin{minipage}{\linewidth}\parindent=\keepfig@pi\relax}{\end{minipage}\par}}{}
\makeatother
\section*{Упражнения}
\addcontentsline{toc}{section}{Упражнения}

% ---------------------------------------------------------------- 4-1
\begin{exercise}[4-1]
Найдите поле разложения многочлена $X^m-1\in\F_p[X]$.
\end{exercise}

\noindent\emph{Решение.} Ответ дан в \rf{2.10c}: если $m=p^kn$, где $p\nmid n$, и $d$~--- порядок $p$ по модулю $n$ (наименьшее
$d\ge1$ с $p^d\equiv1\pmod n$), то поле разложения $X^m-1$ над $\F_p$~--- поле $\F_{p^d}$ из $p^d$ элементов. С точки
зрения теории Галуа (\rf{5.8c}) группа Галуа этого поля порождается автоморфизмом Фробениуса $x\mapsto x^p$, который на
циклической группе $\mu_n$ корней действует возведением в степень $p$; поэтому она отождествляется с подгруппой
$\langle[p]\rangle\subset(\Z/n\Z)^\times$ порядка $d$, а неприводимые множители кругового многочлена $\Phi_n$ по модулю $p$
отвечают смежным классам по этой подгруппе и все имеют степень $d$ (\rf{5.11j}).

Тот же ответ получается из общего описания поля разложения произвольного многочлена над конечным полем.

\begin{proposition}[4.39k]
Пусть $f\in\F_p[X]$~--- непостоянный многочлен \textup(возможно, с кратными множителями\textup), $f_1,\dots,f_r$~--- его
различные унитарные неприводимые множители и $d_i=\deg f_i$. Тогда поле разложения $f$ над $\F_p$ в $\Fpb$ есть $\F_{p^N}$,
где $N=\operatorname{lcm}(d_1,\dots,d_r)$.
\end{proposition}

\begin{proof}
Многочлены $f$ и $f_1\cdots f_r$ имеют одни и те же корни, а значит, одно и то же поле разложения. Многочлен $f_1\cdots f_r$
сепарабелен: каждый $f_i$ имеет $d_i$ различных корней (\rf{4.22}(b)), а различные унитарные неприводимые $f_i$ общих корней не
имеют. По \rf{4.25a} степень его поля разложения над $\F_p$ равна $\operatorname{lcm}(d_1,\dots,d_r)=N$, т.\,е.\ это поле~---
единственное подполе $\F_{p^N}$ поля $\Fpb$ из $p^N$ элементов (\rf{4.23}).
\end{proof}

Для сепарабельного $f$ это последнее утверждение \rf{4.25a}; новое здесь лишь то, что кратные множители не влияют на ответ.
Но существенно, что берутся \emph{различные} неприводимые множители: сам $f$ может иметь кратные множители (например,
$X^m-1$ при $p\mid m$), и тогда он не делит $X^{p^N}-X$~--- этот многочлен не имеет кратных корней. Произведение
$d_1\cdots d_r$ вместо $N$ тоже дало бы поле, содержащее все корни, но, вообще говоря, слишком большое.

Для $f=X^m-1$, $m=p^kn$, имеем $X^m-1=(X^n-1)^{p^k}=\prod_{e\mid n}\Phi_e^{p^k}$; степени неприводимых множителей~--- это
$\operatorname{ord}_e(p)$, $e\mid n$ (\rf{5.11j}), и их НОК равно $\operatorname{ord}_n(p)$ (все они делят $\operatorname{ord}_n(p)$, а
при $e=n$ достигается само $\operatorname{ord}_n(p)$). По \rf{4.39k} это снова ответ \rf{2.10c}. \qed

% ---------------------------------------------------------------- 4-2
\begin{exercise}[4-2]
Найдите группу Галуа многочлена $f=X^4-2X^3-8X-3$ над $\Q$.
\end{exercise}

\noindent\emph{Решение.} Рациональные корни $f$ по \rf{1.11}~--- среди делителей $\pm1,\pm3$ свободного члена. Имеем
$f(1)=-12$, $f(-1)=8$, $f(-3)=156$, но $f(3)=81-54-24-3=0$. Деление даёт
\[
f=(X-3)\,g,\qquad g=X^3+X^2+3X+1 .
\]
Многочлен $g$ неприводим над $\Q$: многочлен степени $3$ приводим ровно тогда, когда у него есть корень, а
$g(1)=6$, $g(-1)=-2$. Так как $3\in\Q$, поля разложения $f$ и $g$ совпадают, и $G_f=G_g$; как группа перестановок четырёх
корней $f$ она оставляет корень $3$ на месте и действует на трёх корнях $g$ транзитивно (\rf{3.27d}). По \rf{4.7} $G_g=A_3$, если
дискриминант $D(g)$~--- квадрат в $\Q$, и $G_g=S_3$ иначе. Сдвиг $X=Y-\frac13$ не меняет разностей корней, а значит, и
дискриминанта, и даёт $g=Y^3+\frac83Y+\frac2{27}$; по формуле для $Y^3+pY+q$ (\rf{4.2e})
\[
D(g)=-4p^3-27q^2=-\frac{4\cdot512}{27}-\frac{27\cdot4}{729}=-\frac{2048+4}{27}=-76 .
\]
Отрицательное число не квадрат в $\Q$, поэтому $G_f=G_g\cong S_3$~--- группа порядка $6$, в $S_4$ это стабилизатор
корня $3$. (По \rf{4.3} знак $D(g)<0$ означает также, что у $g$ один вещественный корень и два комплексно сопряжённых;
комплексное сопряжение и даёт транспозицию в $G_g$.)

\begin{remark}
Транспозицию можно найти и по теореме Дедекинда \rf{4.28}. Из определения дискриминанта
$D(f)=D(g)\prod_{g(\beta)=0}(3-\beta)^2=D(g)\,g(3)^2=-76\cdot46^2=-2^4\cdot19\cdot23^2$, так что $3\nmid D(f)$. По модулю $3$
\[
f\equiv X^4+X^3+X=X(X-1)(X^2-X-1)\pmod 3,
\]
и $X^2-X-1$ неприводим над $\F_3$ (его значения в $0,1,2$ равны $-1,-1,1$). По \rf{4.28} в $G_f$ есть элемент циклового типа
$1+1+2$; поскольку корень $3$ неподвижен, это транспозиция двух корней $g$. Вместе с элементом порядка $3$ (так как
$3\mid|G_g|$) она порождает $S_3$. Заметим, что разложение по модулю $3$ само по себе не показывает приводимость $f$
над $\Q$, а лишь ограничивает возможные разложения: по \rf{1.14} (следствие леммы Гаусса \rf{1.13}) унитарные множители $f$ лежат в $\Z[X]$, и по
модулю $3$ каждый из них~--- произведение некоторых из многочленов $X$, $X-1$, $X^2-X-1$. Множитель $X-3$ приводится в
$X$, а $g$~--- в $(X-1)(X^2-X-1)$.

Так же можно угадать и сам рациональный корень. Если бы $f=q_1q_2$ с унитарными квадратичными $q_1,q_2\in\Z[X]$, то по
модулю $3$ один из них был бы $\equiv X^2-X-1$, а другой $\equiv X(X-1)$ (разложение в $\F_3[X]$ однозначно). По модулю $2$
имеем $f\equiv X^4+1=(X+1)^4$, и по той же однозначности $q_1\equiv q_2\equiv(X+1)^2=X^2+1$. Китайская теорема об остатках,
применённая к каждому коэффициенту, даёт: множитель, сравнимый с $X^2-X$ по модулю $3$ и с $X^2+1$ по модулю $2$, сравним с
$X^2-4X+3=(X-1)(X-3)$ по модулю $6$ (коэффициент при $X$ равен $0$ по модулю $2$ и $-1$ по модулю $3$, т.\,е.\ $-4$ по
модулю $6$; свободный член равен $1$ по модулю $2$ и $0$ по модулю $3$, т.\,е.\ $3$ по модулю $6$), а второй~--- с
$X^2-4X-1$ (их разность $\equiv0\pmod2$ и $\equiv1\pmod3$, т.\,е.\ $\equiv4\pmod6$). Это лишь подсказка, но
она побуждает проверить числа $1$ и $3$, и действительно $f(3)=0$. После этого гипотеза о двух квадратичных множителях
отпадает: $f=(X-3)g$ с неприводимым кубическим $g$, а разложение на неприводимые множители в $\Q[X]$ единственно.
\end{remark}

% ---------------------------------------------------------------- 4-3
\begin{exercise}[4-3]
Найдите степень поля разложения многочлена $X^8-2$ над $\Q$.
\end{exercise}

\noindent\emph{Решение.} Пусть $\theta=\sqrt[8]2>0$ и $\zeta=e^{2\pi i/8}=\frac{1+i}{\sqrt2}$. Корни $X^8-2$~--- это
$\zeta^k\theta$, $k=0,\dots,7$, поэтому поле разложения $L=\Q[\theta,\zeta]$. Так как $\zeta^2=i$ и $\sqrt2=\theta^4$, имеем
$\zeta=(1+i)/\theta^4\in\Q[\theta,i]$ и $i\in\Q[\theta,\zeta]$; значит,
\[
L=\Q[\theta,i].
\]
По критерию Эйзенштейна ($p=2$) многочлен $X^8-2$ неприводим, поэтому $[\Q[\theta]:\Q]=8$. Поле $\Q[\theta]$ лежит в
$\R$ и не содержит $i$, а $i$~--- корень $X^2+1$; значит, $[L:\Q[\theta]]=2$ и
\[
[L:\Q]=16 .
\]

Над $\Q[\theta]$ многочлен раскладывается на неприводимые множители так:
\[
X^8-2=(X-\theta)(X+\theta)(X^2+\theta^2)(X^2-\theta^5X+\theta^2)(X^2+\theta^5X+\theta^2).
\]
Действительно, $X^8-2=(X^4-\theta^4)(X^4+\theta^4)$ и $X^4-\theta^4=(X-\theta)(X+\theta)(X^2+\theta^2)$; множитель же
$X^4+\sqrt2$ над $\Q[\theta]$ ещё раскладывается, потому что $\sqrt2=\theta^4\in\Q[\theta]$ и $2\theta^2=\theta^{10}$:
\[
X^4+\theta^4=(X^2+\theta^2)^2-2\theta^2X^2=(X^2+\theta^2)^2-(\theta^5X)^2 .
\]
Корни квадратичных множителей~--- $\pm i\theta$, $\zeta^{\pm1}\theta$, $\zeta^{\pm3}\theta$ (например,
$\zeta\theta+\zeta^{-1}\theta=\sqrt2\,\theta=\theta^5$ и $\zeta\theta\cdot\zeta^{-1}\theta=\theta^2$); все они невещественны,
так что эти множители неприводимы над $\Q[\theta]\subset\R$. Все корни лежат в $\Q[\theta][i]$, что ещё раз показывает
$[L:\Q[\theta]]=2$.

Ответ $16$, а не $8\cdot\varphi(8)=32$, объясняется тем, что $\Q[\theta]\cap\Q[\zeta]=\Q[\sqrt2]$: элемент
$\sqrt2=\theta^4=\zeta+\zeta^{-1}$ лежит в обоих полях, а $\Q[\zeta]\not\subset\Q[\theta]\subset\R$. В частности,
$[L:\Q[\zeta]]=16/4=4$, а не $8$: над $\Q[\zeta]\ni\sqrt2$ элемент $\theta$~--- корень многочлена $X^4-\sqrt2$. Сейчас мы увидим, что
это соотношение и определяет группу Галуа.

\begin{proposition}[4.39l]
Пусть $G=\Gal(L/\Q)$, $L=\Q[\theta,\zeta]$~--- поле разложения $X^8-2$, и $\chi\colon(\Z/8\Z)^\times\to\{\pm1\}$~--- отображение с
$\chi(\pm1)=1$, $\chi(\pm3)=-1$.
\begin{enumerate}[label=\textup{(\alph*)}]
\item Для каждой пары $a\in\Z/8\Z$, $b\in(\Z/8\Z)^\times$ с $(-1)^a=\chi(b)$ существует ровно один
$\sigma_{a,b}\in G$ с $\sigma_{a,b}(\theta)=\zeta^a\theta$, $\sigma_{a,b}(\zeta)=\zeta^b$, и других элементов в $G$ нет; при
этом $\sigma_{a,b}\sigma_{a',b'}=\sigma_{a+ba',\,bb'}$.
\item Элемент $\tau=\sigma_{1,5}$ имеет порядок $8$, комплексное сопряжение $c=\sigma_{0,-1}$~--- порядок $2$, и
$c\tau c=\tau^3$. Поэтому
\[
G=\langle\tau,c\mid\tau^8=c^2=1,\ c\tau c^{-1}=\tau^3\rangle=\mathrm{SD}_{16}
\]
~--- полудиэдральная \textup(квазидиэдральная\textup) группа порядка $16$; $\langle\tau\rangle=\Gal(L/\Q[i])$.
\item В $G$ один элемент порядка $1$, пять элементов порядка $2$, шесть порядка $4$ и четыре порядка $8$. Центр $G$
равен $\langle\tau^4\rangle$, и $G/\langle\tau^4\rangle\cong D_4$.
\end{enumerate}
\end{proposition}

\begin{proof}
(a) Элемент $\sigma\in G$ переводит $\theta$ в корень $X^8-2$, т.\,е.\ в $\zeta^a\theta$, а $\zeta$~--- в корень многочлена
$\Phi_8=X^4+1$, т.\,е.\ в первообразный корень $\zeta^b$, $b$ нечётно; так как $L=\Q[\theta,\zeta]$, пара $(a,b)$ определяет
$\sigma$. Применим $\sigma$ к равенству $\theta^4=\sqrt2=\zeta+\zeta^{-1}$. Слева получим $\zeta^{4a}\theta^4=(-1)^a\sqrt2$,
справа $\zeta^b+\zeta^{-b}$, что равно $\sqrt2$ при $b=\pm1$ и $\zeta^3+\zeta^{-3}=-\sqrt2$ при $b=\pm3$. Значит,
$(-1)^a=\chi(b)$. Таких пар ровно $8\cdot4/2=16$ (при каждом $b$ годятся четыре значения $a$ нужной чётности), а $|G|=[L:\Q]=16$
(\rf{3.10}). Поэтому каждая такая пара действительно отвечает элементу $G$. Наконец,
$\sigma_{a,b}(\sigma_{a',b'}(\theta))=\sigma_{a,b}(\zeta^{a'}\theta)=\zeta^{ba'}\zeta^{a}\theta$ и
$\sigma_{a,b}(\sigma_{a',b'}(\zeta))=\zeta^{bb'}$.

(b) Пары $(1,5)$ и $(0,-1)$ допустимы: $(-1)^1=\chi(5)$, $(-1)^0=\chi(-1)$. По правилу умножения $\tau^2=\sigma_{6,1}$,
$\tau^3=\sigma_{7,5}$, $\tau^4=\sigma_{4,1}\ne1$, $\tau^8=\sigma_{0,1}=1$; значит, порядок $\tau$ равен $8$. Отображение $c$
оставляет $\theta\in\R$ на месте и переводит $\zeta$ в $\bar\zeta=\zeta^{-1}$, т.\,е.\ это комплексное сопряжение. Далее
$\tau c=\sigma_{1,3}$ и $c\tau c=\sigma_{0,-1}\sigma_{1,3}=\sigma_{-1,-3}=\sigma_{7,5}=\tau^3$. Все $\tau^k$ оставляют на месте
$i=\zeta^2$ (так как $5\cdot2\equiv2\pmod8$), а $c(i)=-i$; поэтому $c\notin\langle\tau\rangle$, $\langle\tau\rangle$~---
подгруппа индекса $2$, т.\,е.\ $\langle\tau\rangle=\Gal(L/\Q[i])$, и $G=\langle\tau\rangle\sqcup\langle\tau\rangle c$.
Соотношения $\tau^8=c^2=1$, $c\tau c=\tau^3$ позволяют записать любое произведение образующих в виде $\tau^k$ или $\tau^kc$,
так что группа с этими образующими и соотношениями имеет не больше $16$ элементов; $G$ удовлетворяет им и имеет
порядок $16$, поэтому $G$ и есть эта группа.

(c) Порядок $\tau^k$ равен $8/\gcd(k,8)$: четыре элемента $\tau,\tau^3,\tau^5,\tau^7$ порядка $8$, два элемента
$\tau^2,\tau^6$ порядка $4$ и $\tau^4$ порядка $2$. Далее $(\tau^kc)^2=\tau^k(c\tau^kc)=\tau^{k}\tau^{3k}=\tau^{4k}$; это $1$
при чётном $k$ (элементы $c,\tau^2c,\tau^4c,\tau^6c$ порядка $2$) и $\tau^4\ne1$ при нечётном $k$ (элементы
$\tau c,\tau^3c,\tau^5c,\tau^7c$ порядка $4$). Элемент $\tau^kc$ не коммутирует с $\tau$ ($\tau\cdot\tau^kc=\tau^{k+1}c$, а
$\tau^kc\cdot\tau=\tau^{k+3}c$), а $\tau^k$ коммутирует с $c$ ровно при $\tau^{3k}=\tau^k$, т.\,е.\ $k\in\{0,4\}$. Значит, центр
равен $\langle\tau^4\rangle$. В факторгруппе по нему образ $\bar\tau$ имеет порядок $4$ и $\bar c\bar\tau\bar c=\bar\tau^3=\bar\tau^{-1}$,
т.\,е.\ это группа диэдра $D_4$ порядка $8$.
\end{proof}

\begin{remark}
Сопоставление $\sigma_{a,b}\mapsto b$~--- это ограничение на $\Q[\zeta]$ (так как $\sigma_{a,b}(\zeta)=\zeta^b$), т.\,е.\
сюръективный гомоморфизм
\[
G\to\Gal(\Q[\zeta]/\Q)\cong(\Z/8\Z)^\times=\{\pm1\}\times\{1,5\}\cong(\Z/2\Z)^2
\]
(\rf{5.11a}, \rf{5.11g}); его ядро~--- $\Gal(L/\Q[\zeta])=\{\sigma_{a,1}\mid a\text{ чётно}\}=\langle\tau^2\rangle$ порядка $4$. Правило
умножения из (a) показывает, что $G$~--- подгруппа индекса $2$ в полупрямом произведении $\Z/8\Z\rtimes(\Z/8\Z)^\times$
(пары $(a,b)$ с умножением $(a,b)(a',b')=(a+ba',bb')$; это группа преобразований $x\mapsto bx+a$ кольца $\Z/8\Z$). Вообще,
для $X^n-c\in\Q[X]$, $c\ne0$, с корнем $\theta$ и первообразным корнем $\zeta$ степени $n$ из $1$ то же рассуждение даёт
вложение группы Галуа в $\Z/n\Z\rtimes(\Z/n\Z)^\times$, $\sigma\mapsto(a,b)$, где $\sigma(\theta)=\zeta^a\theta$,
$\sigma(\zeta)=\zeta^b$; оно является изоморфизмом ровно тогда, когда степень поля разложения равна $n\varphi(n)$. Здесь
легко ошибиться: <<естественное>> отображение $\theta\mapsto\zeta\theta$, $\zeta\mapsto\zeta$ (пара $(1,1)$) автоморфизмом
\emph{не} является~--- оно переводило бы $\sqrt2=\theta^4$ в $\zeta^4\theta^4=-\sqrt2$, а $\sqrt2=\zeta+\zeta^{-1}$ оставляло бы
на месте. Именно соотношение $(-1)^a=\chi(b)$ отсекает половину пар.
\end{remark}

\noindent\emph{Почему это именно $\mathrm{SD}_{16}$.} Неабелевых групп порядка $16$ с элементом $\tau$ порядка $8$ четыре (это
известный факт теории групп): группа диэдра $D_8$ порядка $16$, так называемая модулярная группа $M_{16}$ и полудиэдральная
группа $\mathrm{SD}_{16}$~--- в них есть элемент $c$ порядка $2$ с $c\tau c^{-1}=\tau^k$, где $k=7,5,3$ соответственно,~---
а также обобщённая группа кватернионов $Q_{16}$ (в ней $c\tau c^{-1}=\tau^{-1}$, $c^2=\tau^4$). Их различает число элементов порядка $2$: в $D_8$ их девять, в
$M_{16}$~--- три ($\tau^4$, $c$, $\tau^4c$), в $Q_{16}$~--- один, в $\mathrm{SD}_{16}$~--- пять. В нашей группе их пять, так что это
$\mathrm{SD}_{16}$. (Вычисление в PARI: \texttt{galoisidentify} для поля разложения $X^8-2$ даёт группу $[16,8]$ из таблицы
малых групп, т.\,е.\ $\mathrm{SD}_{16}$.)

\subsection*{Подполя поля разложения $X^8-2$}

У $\mathrm{SD}_{16}$ ровно $15$ подгрупп: $\{1\}$; пять подгрупп порядка $2$ (по одной на элемент порядка $2$); три
циклические подгруппы порядка $4$ (шесть элементов порядка $4$, по два в каждой) и две подгруппы $\{1,\tau^4,s,\tau^4s\}$,
$s\in\{c,\tau^2c\}$, изоморфные $V_4$: элементы $\tau^{2j}c$ и $\tau^{2k}c$ коммутируют ровно тогда, когда $\tau^{2j}c\,\tau^{2k}c=\tau^{2j-2k}$ имеет
порядок $\le2$, т.\,е.\ при $j-k\equiv0,2\pmod4$; три подгруппы порядка $8$ (подгруппы индекса $2$ содержат все квадраты,
т.\,е.\ $\langle\tau^2\rangle$, и отвечают трём подгруппам порядка $2$ в $G/\langle\tau^2\rangle\cong(\Z/2\Z)^2$); и сама $G$.
По основной теореме \rf{3.17} подполей $L$ тоже $15$. Чтобы найти неподвижное поле $L^H$, достаточно указать элемент $x\in L$,
неподвижный относительно $H$, степень которого равна индексу $(G:H)$: тогда $\Q[x]\subset L^H$ и
$[\Q[x]:\Q]=(G:H)=[L^H:\Q]$, т.\,е.\ $\Q[x]=L^H$. Результат~--- в таблице.

\begin{center}
\begin{keepfig}
\centering
\small
\begin{tabular}{llll}
\toprule
$H$ & $|H|$, строение & $L^H=\Q[x]$ & пояснение \\
\midrule
$G$ & $16$ & $\Q$ & \\
$\langle\tau\rangle$ & $8$, $C_8$ & $\Q[i]$ & $5\cdot2\equiv2\pmod 8$ \\
$\langle\tau^2,c\rangle$ & $8$, $D_4$ & $\Q[\sqrt2]$ & $\sqrt2=\theta^4\mapsto\zeta^{4a}\theta^4$, $a$ чётно \\
$\langle\tau^2,\tau c\rangle$ & $8$, $Q_8$ & $\Q[\sqrt{-2}]$ & $\tau c\colon\theta\mapsto\zeta\theta,\ i\mapsto-i$ \\
$\langle\tau^2\rangle$ & $4$, $C_4$ & $\Q[\zeta]=\Q[i,\sqrt2]$ & $\tau^2=\sigma_{6,1}$ \\
$\langle\tau^4,c\rangle$ & $4$, $V_4$ & $\Q[\theta^2]=\Q[\sqrt[4]2]$ & $\theta^2$~--- корень $X^4-2$ \\
$\langle\tau^4,\tau^2c\rangle$ & $4$, $V_4$ & $\Q[i\theta^2]$ & $i\theta^2$~--- корень $X^4-2$ \\
$\langle\tau c\rangle$ & $4$, $C_4$ & $\Q[(1+i)\theta^2]$ & корень $X^4+8$ \\
$\langle\tau^3c\rangle$ & $4$, $C_4$ & $\Q[(1-i)\theta^2]$ & корень $X^4+8$ \\
$\langle\tau^4\rangle$ & $2$ & $\Q[\theta^2,i]$ & поле разложения $X^4-2$ \\
$\langle c\rangle$ & $2$ & $\Q[\theta]$ & корень $X^8-2$ \\
$\langle\tau^4c\rangle$ & $2$ & $\Q[i\theta]$ & корень $X^8-2$ \\
$\langle\tau^6c\rangle$ & $2$ & $\Q[\zeta\theta]$ & корень $X^8-2$ \\
$\langle\tau^2c\rangle$ & $2$ & $\Q[\zeta^3\theta]$ & корень $X^8-2$ \\
$\{1\}$ & $1$ & $L$ & \\
\bottomrule
\end{tabular}
\end{keepfig}
\end{center}

Проверки неподвижности прямые. Элементы порядка $2$ вне центра~--- это $\tau^{2k}c=\sigma_{2k,-1}$ (по правилу умножения
$\tau^{2k}=\sigma_{a,1}$ с $a\equiv-2k\pmod 8$, так что множество $\{\tau^{2k}c\}$~--- это все $\sigma_{2k,-1}$), и
$\sigma_{2k,-1}(\zeta^k\theta)=\zeta^{-k}\zeta^{2k}\theta=\zeta^k\theta$; при этом $\tau^4c=\sigma_{4,-1}$, $\tau^6c=\sigma_{2,-1}$,
$\tau^2c=\sigma_{6,-1}$, что и даёт четыре строки с полями $\Q[\zeta^k\theta]$, $k=0,1,2,3$ (поле $\Q[\zeta^{k+4}\theta]=\Q[-\zeta^k\theta]$
то же самое). Для $\tau c=\sigma_{1,3}$ имеем $\theta^2\mapsto\zeta^2\theta^2=i\theta^2$, $i\mapsto\zeta^6=-i$, поэтому
$(1+i)\theta^2\mapsto(1-i)\,i\theta^2=(1+i)\theta^2$; аналогично $\tau^3c=\sigma_{7,3}$ переводит $\theta^2$ в $\zeta^{14}\theta^2=-i\theta^2$
и оставляет на месте $(1-i)\theta^2$. Число $x=(1\pm i)\theta^2$ удовлетворяет $x^2=\pm2i\sqrt2$, $x^4=-8$, а многочлен $X^4+8$
неприводим над $\Q$: рациональных корней у него нет, а из $X^4+8=(X^2+aX+b)(X^2-aX+c)$ с $a,b,c\in\Z$ (\rf{1.14}) следует
$a(c-b)=0$, $b+c=a^2$, $bc=8$, т.\,е.\ либо $a=0$, $c=-b$, $-b^2=8$, либо $b=c$, $b^2=8$~--- невозможно. Для $\tau^2c=\sigma_{6,-1}$:
$\theta^2\mapsto\zeta^{12}\theta^2=-\theta^2$, $i\mapsto-i$, так что $i\theta^2\mapsto i\theta^2$; а $(i\theta^2)^4=2$.

\begin{center}
\begin{keepfig}
\centering
\begin{tikzpicture}[every node/.style={font=\small},x=1cm,y=1cm]
\node[cnew] (Q) at (0,0) {$\Q$};
\node[cnew] (s2) at (-3.4,1.7) {$\Q[\sqrt2]$};
\node[cnew] (i) at (0.6,1.7) {$\Q[i]$};
\node[cnew] (sm2) at (4.2,1.7) {$\Q[\sqrt{-2}]$};
\node (t2) at (-5.4,3.4) {$\Q[\theta^2]$};
\node (it2) at (-2.4,3.4) {$\Q[i\theta^2]$};
\node[cnew] (z) at (0.6,3.4) {$\Q[\zeta]$};
\node (pt2) at (3.1,3.4) {$\Q[(1{+}i)\theta^2]$};
\node (mt2) at (5.9,3.4) {$\Q[(1{-}i)\theta^2]$};
\node (t) at (-6.4,5.1) {$\Q[\theta]$};
\node (it) at (-4.6,5.1) {$\Q[i\theta]$};
\node (zt) at (-2.8,5.1) {$\Q[\zeta\theta]$};
\node (z3t) at (-1.0,5.1) {$\Q[\zeta^3\theta]$};
\node[cnew] (t2i) at (2.6,5.1) {$\Q[\theta^2,i]$};
\node[cnew] (L) at (-1.4,6.8) {$L=\Q[\theta,i]$};
\draw (Q)--(s2); \draw (Q)--(i); \draw (Q)--(sm2);
\draw (s2)--(t2); \draw (s2)--(it2); \draw (s2)--(z);
\draw (i)--(z);
\draw (sm2)--(z); \draw (sm2)--(pt2); \draw (sm2)--(mt2);
\draw (t2)--(t); \draw (t2)--(it); \draw (t2)--(t2i);
\draw (it2)--(zt); \draw (it2)--(z3t); \draw (it2)--(t2i);
\draw (z)--(t2i); \draw (pt2)--(t2i); \draw (mt2)--(t2i);
\foreach \n in {t,it,zt,z3t,t2i} \draw (L)--(\n);
\node[left] at (-7.0,0) {$1$};
\node[left] at (-7.0,1.7) {$2$};
\node[left] at (-7.0,3.4) {$4$};
\node[left] at (-7.0,5.1) {$8$};
\node[left] at (-7.0,6.8) {$16$};
\end{tikzpicture}
\figcap{Все $15$ подполей поля разложения $L$ многочлена $X^8-2$ ($\theta=\sqrt[8]2$, $\zeta=e^{2\pi i/8}$); слева~--- степень
над $\Q$. Красным выделены подполя, являющиеся расширениями Галуа поля $\Q$: им отвечают нормальные подгруппы $\{1\}$,
$\langle\tau^4\rangle$, $\langle\tau^2\rangle$, три подгруппы порядка $8$ и $G$. Поле $\Q[\theta]$ содержит единственную цепочку
$\Q\subset\Q[\sqrt2]\subset\Q[\sqrt[4]2]\subset\Q[\theta]$.}
\end{keepfig}
\end{center}

\begin{remark}
\emph{Подполя поля разложения $X^4-2$.} Поле $\Q[\theta^2,i]=\Q[\sqrt[4]2,i]$~--- поле разложения $X^4-2$ (\rf{4.13}), и ему отвечает центр $\langle\tau^4\rangle$.
Поэтому его подполя~--- это десять полей нижней части рисунка (те, что лежат под $\Q[\theta^2,i]$), а группа
$\Gal(\Q[\sqrt[4]2,i]/\Q)\cong G/\langle\tau^4\rangle\cong D_4$ (\rf{3.17}). Запишем её явно как группу перестановок корней
\[
1\colon\sqrt[4]2,\qquad 2\colon i\sqrt[4]2,\qquad 3\colon-\sqrt[4]2,\qquad 4\colon -i\sqrt[4]2
\]
(вершин квадрата на комплексной плоскости). Ограничение $r$ элемента $\tau$ оставляет $i$ на месте и переводит
$\sqrt[4]2=\theta^2$ в $\zeta^2\theta^2=i\sqrt[4]2$, а потому $i\sqrt[4]2\mapsto i\cdot i\sqrt[4]2=-\sqrt[4]2$ и т.\,д.: это поворот
$r=(1\,2\,3\,4)$ на $90^\circ$ порядка $4$. Комплексное сопряжение~--- отражение $(2\,4)$. Всего
\[
D_4=\{1,\ (1\,2\,3\,4),\ (1\,3)(2\,4),\ (1\,4\,3\,2),\ (2\,4),\ (1\,3),\ (1\,2)(3\,4),\ (1\,4)(2\,3)\},
\]
и соответствие подгрупп и полей таково:
\[
\begin{gathered}
\langle r\rangle\leftrightarrow\Q[i],\qquad \{1,(1\,3),(2\,4),(1\,3)(2\,4)\}\leftrightarrow\Q[\sqrt2],\\
\{1,(1\,2)(3\,4),(1\,3)(2\,4),(1\,4)(2\,3)\}\leftrightarrow\Q[\sqrt{-2}],\qquad
\langle r^2\rangle\leftrightarrow\Q[i,\sqrt2],\\
\langle(2\,4)\rangle\leftrightarrow\Q[\sqrt[4]2],\qquad
\langle(1\,3)\rangle\leftrightarrow\Q[i\sqrt[4]2],\\
\langle(1\,2)(3\,4)\rangle\leftrightarrow\Q[(1+i)\sqrt[4]2],\qquad \langle(1\,4)(2\,3)\rangle\leftrightarrow\Q[(1-i)\sqrt[4]2].
\end{gathered}
\]
Например, $(1\,2)(3\,4)$ меняет местами $\sqrt[4]2$ и $i\sqrt[4]2$ (это ограничение $\tau c$), поэтому оставляет на месте их
сумму $(1+i)\sqrt[4]2$, а её степень равна $4$ (корень $X^4+8$). Четыре подгруппы порядка $2$, порождённые отражениями, не
нормальны: например, $r(2\,4)r^{-1}=(1\,3)$, и соответственно поля $\Q[\sqrt[4]2]$ и $\Q[i\sqrt[4]2]$ сопряжены, а
$\Q[(1+i)\sqrt[4]2]$ и $\Q[(1-i)\sqrt[4]2]$ сопряжены между собой; остальные подгруппы нормальны. Так как
$(1\,2)(3\,4)\cdot(1\,3)=(1\,4\,3\,2)$, произведение двух отражений~--- поворот. Наконец, кубическая резольвента
$X^4-2$ равна $X^3+8X=X(X^2+8)$ (\rf{4.9}), её поле разложения $\Q[\sqrt{-2}]$ отвечает подгруппе
$G\cap V=\{1,(1\,2)(3\,4),(1\,3)(2\,4),(1\,4)(2\,3)\}$, как и должно быть по \rf{4.8}.
\end{remark}

% ---------------------------------------------------------------- 4-4
\begin{exercise}[4-4]
Приведите пример расширения полей $E/F$ степени $4$, для которого не существует поля $M$ с $F\subset M\subset E$,
$[M:F]=2$.
\end{exercise}

\noindent\emph{Что не подходит.} (1) Расширения Галуа: группа $G=\Gal(E/F)$ порядка $4$ по теореме Коши содержит подгруппу
$H$ порядка $2$, и $M=E^H$ имеет степень $(G:H)=2$ (\rf{3.17}). В частности, не подходят расширения конечных полей: они
являются расширениями Галуа (а $\F_{q^4}\supset\F_{q^2}\supset\F_q$ видно и по \rf{4.23}). (2) Поле $\Q[\sqrt[4]2]$: оно
содержит $\Q[\sqrt2]$. По предыдущему замечанию то же верно для $\Q[i\sqrt[4]2]\supset\Q[\sqrt2]$ и для
$\Q[(1\pm i)\sqrt[4]2]\supset\Q[\sqrt{-2}]$. Как мы сейчас увидим, причина в том, что группа Галуа $X^4-2$~--- это $D_4$.
(3) Несепарабельные расширения степени $4$ (см.\ замечание ниже).

\medskip\noindent\emph{Решение.} Пусть $f=X^4+8X+12$, $\alpha$~--- его корень, $E=\Q[\alpha]$, $E_f$~--- поле разложения $f$.

\emph{$f$ неприводим и $G_f=A_4$.} Это пример \rf{4.13a}: у $f$ нет вещественных корней, разложению на квадратичные множители
отвечал бы корень вида $a^2$ многочлена $X^3-48X-64$ (\rf{4.9i}), которого нет, так что $f$ неприводим; его кубическая
резольвента $g=X^3-48X-64$ неприводима, а $D(g)=331776=576^2$~--- квадрат; по \rf{4.9d} $G_f=A_4$, $|G_f|=12$. В частности,
$[E:\Q]=4$.

\emph{Промежуточных полей нет.} Пусть $\Q\subset M\subset E$, $[M:\Q]=2$. По \rf{3.17} подгруппа $H=\Gal(E_f/M)$ имеет в
$G_f=A_4$ индекс $[M:\Q]=2$. Но таких подгрупп нет:

\begin{lemma}[4.39m]
При $n\ge1$ в знакопеременной группе $A_n$ нет подгрупп индекса $2$.
\end{lemma}

\begin{proof}
Подгруппа $H$ индекса $2$ в группе $K$ нормальна, и $K/H$~--- группа порядка $2$; поэтому $k^2\in H$ для всех $k\in K$.
Каждый тройной цикл $\rho$~--- квадрат: $\rho=\rho^4=(\rho^2)^2$. Значит, подгруппа индекса $2$ в $A_n$ содержала бы все
тройные циклы, а при $n\ge3$ они порождают $A_n$; при $n\le2$ группа $A_n$ тривиальна. В обоих случаях $H=A_n$~---
противоречие. (Для $A_4$ можно и сосчитать: восемь тройных циклов не помещаются в подгруппу порядка $6$.)
\end{proof}

Итак, $E=\Q[\alpha]$, $\alpha^4+8\alpha+12=0$,~--- искомый пример. Более того, между $\Q$ и $E$ нет вообще никаких
промежуточных полей: по \rf{1.20} их степень могла бы быть только $2$.

\medskip\noindent\emph{Второй пример: группа $S_4$.} Пусть $f$~--- неприводимая квартика с $G_f=S_4$ (например,
$X^4-4X+2$, \rf{4.10}), $E=\Q[\alpha]$. Промежуточные поля $\Q\subset M\subset E$ отвечают подгруппам $K$ с
$\Gal(E_f/E)\subset K\subset S_4$, а $\Gal(E_f/E)$~--- стабилизатор корня $\alpha$, т.\,е.\ $S_3$. Если $S_3\subsetneq K$, то $K$
транзитивна (она переводит $\alpha$ в другой корень, а $S_3$ транзитивно переставляет три остальных) и содержит
транспозицию и тройной цикл; по \rf{4.31} $K=S_4$. Значит, промежуточных полей снова нет. То же рассуждение показывает: если
$G_f=S_n$, то в $F[\alpha]$ нет промежуточных полей, а если $G_f=A_n$, то в $F[\alpha]$ нет подполей степени $2$ (\rf{4.39m}).

\medskip
Какие квартики дают пример, а какие нет, описывает следующее предложение.

\begin{proposition}[4.39n]
Пусть $\operatorname{char}F\ne2$, $f\in F[X]$~--- неприводимый сепарабельный многочлен степени $4$ с корнями
$\alpha_1,\dots,\alpha_4$ в поле разложения $E_f$, $E=F[\alpha_1]$. Следующие условия равносильны:
\begin{enumerate}[label=\textup{(\alph*)}]
\item существует поле $M$, $F\subsetneq M\subsetneq E$ \textup(тогда $[M:F]=2$\textup);
\item группа $G_f$ сохраняет разбиение корней на две пары \textup(т.\,е.\ содержится в одной из трёх подгрупп $D_4$
группы $S_4$, \rf{4.7j}(c)\textup);
\item кубическая резольвента $g$ многочлена $f$ \textup(\rf{4.8a}, \rf{4.9}\textup) имеет корень в $F$.
\end{enumerate}
В частности, примеры к упражнению 4-4 среди простых расширений $F[\alpha]$ степени $4$~--- это в точности случаи
$G_f=A_4$ и $G_f=S_4$.
\end{proposition}

\begin{proof}
Пусть $G=G_f$, $G_1=\Gal(E_f/E)$~--- стабилизатор $\alpha_1$; так как $G$ транзитивна (\rf{3.27d}), $(G:G_1)=4$.

(a)$\Rightarrow$(b). Пусть $H=\Gal(E_f/M)$; по \rf{3.17} $G_1\subset H$ и $(G:H)=2$. Подгруппа $H$ нормальна, поэтому $\rho\in G$
переводит $H$-орбиты в $H$-орбиты ($\rho(H\beta)=H(\rho\beta)$), и так как $G$ транзитивна, все $H$-орбиты одного размера.
Орбита $\alpha_1$ имеет размер $|H|/|H\cap G_1|=|H|/|G_1|=2$. Значит, $H$-орбиты разбивают корни на две пары, и $G$ это
разбиение сохраняет.

(b)$\Rightarrow$(a). Пусть $G$ сохраняет разбиение $\{\{\alpha_1,\alpha_2\},\{\alpha_3,\alpha_4\}\}$, и
$K=\{\rho\in G\mid \rho\{\alpha_1,\alpha_2\}=\{\alpha_1,\alpha_2\}\}$. Элемент, оставляющий $\alpha_1$ на месте, переводит пару,
содержащую $\alpha_1$, в себя; значит, $G_1\subset K$. Так как $G$ транзитивна, некоторый $\rho\in G$ переводит $\alpha_1$ в
$\alpha_3$ и меняет пары местами; поэтому $K\ne G$, а $K$~--- ядро гомоморфизма $G\to\Sym\bigl(\{\{\alpha_1,\alpha_2\},\{\alpha_3,\alpha_4\}\}\bigr)$, так что $(G:K)=2$. Поле $M=E_f^K$ лежит в $E_f^{G_1}=E$ и имеет степень $2$ над $F$.
(Явно: $M=F[\alpha_1\alpha_2]$, если $\alpha_1\alpha_2\ne\alpha_3\alpha_4$.)

(b)$\Leftrightarrow$(c). Корни резольвенты~--- $\theta_1=\alpha_1\alpha_2+\alpha_3\alpha_4$, $\theta_2=\alpha_1\alpha_3+\alpha_2\alpha_4$,
$\theta_3=\alpha_1\alpha_4+\alpha_2\alpha_3$ (\rf{4.8a}); они попарно различны (\rf{4.7j}(a)), и $\sigma\in G$ переводит $\theta_i$ в
$\theta_{\pi(\sigma)(i)}$ (\rf{4.7j}(b)). Поэтому $\sigma$ оставляет $\theta_1$ на месте ровно тогда, когда сохраняет разбиение
$\{\{\alpha_1,\alpha_2\},\{\alpha_3,\alpha_4\}\}$, т.\,е.\ лежит в $D^{(1)}$ (\rf{4.7j}(c)); то же для $\theta_2,\theta_3$ и двух других
разбиений. По \rf{3.17} $\theta_i\in F\iff G$ оставляет $\theta_i$ на месте. Значит, резольвента имеет корень в $F$ $\iff$ $G$
сохраняет одно из разбиений (ср.\ \rf{4.9e}).

Последнее утверждение: по \rf{4.9d} резольвента имеет корень в $F$ ровно при $G_f\in\{V_4,C_4,D_4\}$; любое сепарабельное
расширение степени $4$ просто (теорема о примитивном элементе \rf{5.1}).
\end{proof}

Для $X^4-2$ корень $0$ резольвенты $X^3+8X$ отвечает разбиению корней на пары $\{\sqrt[4]2,-\sqrt[4]2\}$,
$\{i\sqrt[4]2,-i\sqrt[4]2\}$, и $M=\Q[\alpha_1\alpha_2]=\Q[\sqrt2]$ ($\alpha_1\alpha_2=-\sqrt2$)~--- как и было найдено выше.

\begin{remark}
\emph{Несепарабельные расширения степени $4$ не дают примеров.} Пусть $\operatorname{char}F=p>0$, $[E:F]=4$, и $E/F$
несепарабельно. Пусть $E_{\mathrm{sep}}$~--- сепарабельное замыкание $F$ в $E$; по \rf{3.15a}(d) $[E:E_{\mathrm{sep}}]$~--- степень
$p$, большая $1$ и делящая $4$, так что $p=2$. Если $[E_{\mathrm{sep}}:F]=2$, то $M=E_{\mathrm{sep}}$. Если $E_{\mathrm{sep}}=F$, то
$E/F$ чисто несепарабельно, и минимальный многочлен любого $x\in E\smallsetminus F$ имеет вид $X^{2^e}-a$ (\rf{3.15a}(b)), $e\in\{1,2\}$.
Если для какого-то $x$ $e=1$, то $M=F[x]$; иначе $E=F[x]$, $x^4=a$, и $M=F[x^2]$ имеет степень $2$.
\end{remark}

% ---------------------------------------------------------------- 4-5
\begin{exercise}[4-5]
Перечислите все неприводимые многочлены степени $3$ над $\F_7$ за $10$ секунд или быстрее (их $112$).
\end{exercise}

\noindent\emph{Решение.} \emph{Число.} По \rf{4.22b} многочлен $X^{7^3}-X$~--- произведение всех унитарных неприводимых
многочленов над $\F_7$, степень которых делит $3$, причём каждый входит один раз. Неприводимых многочленов степени $1$ семь,
поэтому для числа $N$ неприводимых кубических $343=7+3N$, т.\,е.\ $N=112$ (это \rf{4.22c}(a): $N_7(3)=(7^3-7)/3$). (Можно и прямо: из $343$ унитарных кубических
многочленов приводимы те, что имеют корень: $\binom{9}{3}=84$ произведений трёх линейных множителей и
$7\cdot\frac{49-7}2=147$ произведений линейного на неприводимый квадратичный; $343-84-147=112$.) На компьютере ответ даёт
одна команда: разложение $X^{343}-X$ по модулю $7$ (в PARI \texttt{factormod(x\^{}343-x,7)}) содержит $7$ линейных и
$112$ кубических множителей.

\emph{Быстрое перечисление вручную.} Кубический многочлен неприводим тогда и только тогда, когда у него нет корня в $\F_7$.
Используем две замены, сохраняющие неприводимость.

(1) \emph{Сдвиги.} Для $h=X^3+aX^2+bX+c$ коэффициент при $X^2$ у $h(X+t)$ равен $a+3t$; так как $3$ обратимо в $\F_7$, среди
семи сдвигов $h(X+t)$, $t\in\F_7$, ровно один имеет вид $X^3+pX+q$. Поэтому достаточно найти $112/7=16$ неприводимых
многочленов вида $X^3+pX+q$; все $112$~--- это их сдвиги.

(2) \emph{Растяжения.} Для $u\in\F_7^\times$ многочлен $u^{-3}(X^3+pX+q)\big|_{X\to uX}=X^3+u^{-2}pX+u^{-3}q$ неприводим
вместе с исходным. В $\F_7$ квадраты~--- это $1,2,4$, кубы~--- это $\pm1$, и $u$ определяется парой $(u^2,u^3)$
($u=u^3/u^2$).

Теперь: $X^3+q$ неприводим $\iff$ $-q$ не куб $\iff$ $q\in\{2,3,4,5\}$ (четыре многочлена). Многочлен $X^3+X+1$ не имеет корней
(значения в $0,1,\dots,6$: $1,3,4,3,6,5,6$), и его растяжения дают шесть многочленов $X^3+pX\pm1$, $p\in\{1,2,4\}$. Многочлен
$X^3+3X+2$ не имеет корней (значения $2,6,2,3,1,2,5$), и его растяжения дают шесть многочленов $X^3+pX\pm2$, $p\in\{3,5,6\}$.
Получено $4+6+6=16$ различных многочленов~--- значит, это все. Итак, неприводимые кубические многочлены над $\F_7$~--- это
\[
(X+t)^3+q\ \ (q=2,3,4,5),\quad (X+t)^3+p(X+t)\pm1\ \ (p=1,2,4),\quad (X+t)^3+p(X+t)\pm2\ \ (p=3,5,6),
\]
$t\in\F_7$: $16\cdot7=112$ штук (список проверен на компьютере перебором). Отметим, что дискриминант
$-4p^3-27q^2$ у всех шестнадцати~--- ненулевой квадрат в $\F_7$, как и должно быть: группа Галуа неприводимого многочлена над
конечным полем циклическая (\rf{4.20}), для кубического это $A_3$ (\rf{4.2}). Обратное неверно: у $X^3-X=X(X-1)(X+1)$ дискриминант
$4$~--- тоже квадрат. \qed

% ---------------------------------------------------------------- 4-6
\begin{exercise}[4-6]
<<Это наводящий на размышления вопрос: немногие аспиранты знали бы, как подступиться к задаче определения группы Галуа,
скажем, многочлена
\[
X^6+2X^5+3X^4+4X^3+5X^2+6X+7\text{>>}
\]
\textup(над $\Q$\textup). (a) Можете ли вы её найти? (b) Можете ли вы найти её, не пользуясь командой \texttt{polgalois}
в PARI?
\end{exercise}

\noindent\emph{Ответ.} $G_f\cong\mathrm{PGL}_2(\F_5)\cong S_5$~--- группа порядка $120$ обратимых матриц $2\times2$ над $\F_5$
по модулю скалярных, действующая на шести корнях так же транзитивно, как $\mathrm{PGL}_2(\F_5)$ действует на шести точках
проективной прямой $\F_5\cup\{\infty\}$ дробно-линейными преобразованиями $z\mapsto\frac{az+b}{cz+d}$. (a) Команда
\texttt{polgalois} в PARI выдаёт \texttt{[120, -1, 14, "L(6):2 = PGL(2,5) = S\_5(6)"]}: порядок $120$, группа не лежит
в $A_6$ (знак $-1$), транзитивная группа степени $6$ номер $14$ в стандартной таблице (вычислено на компьютере). Обоснуем ответ.

\medskip\noindent\emph{Решение} (b). Обозначим многочлен через $f$. Полезно заметить, что $(X-1)^2f=X^8-8X+7$, так что
корни $f$~--- это корни трёхчлена $X^8-8X+7$, отличные от двойного корня $1$. Дискриминант $f$ равен $-2^{16}\cdot7^4$
(вычислено, например, через результант, \rf{4.36}). Он отрицателен, т.\,е.\ не квадрат, и по \rf{4.2} $G_f\not\subset A_6$.

\emph{Шаг 1: теорема Дедекинда.} Для простых $p\ne2,7$ (не делящих дискриминант) разложение $f$ по модулю $p$ даёт цикловой
тип некоторого элемента $G_f$ (\rf{4.28}). Вычисление даёт:
\begin{center}
\begin{tabular}{lcccccc}
\toprule
$p$ & $3$ & $5$ & $13$ & $19$ & $61$ & $79$\\
степени множителей $f\bmod p$ & $6$ & $3+3$ & $1+5$ & $1+1+4$ & $1+1+2+2$ & $2+2+2$\\
\bottomrule
\end{tabular}
\end{center}
В частности, $f$ неприводим (уже по модулю $3$), $G_f$ транзитивна, а порядок $G_f$ делится на $3$, $4$ и $5$, т.\,е. на $60$.

\emph{Шаг 2: кандидаты.} Транзитивные подгруппы $S_6$ порядка, делящегося на $60$,~--- это $S_6$, $A_6$ и <<экзотические>> копии
$\mathrm{PGL}_2(\F_5)\cong S_5$ и $\mathrm{PSL}_2(\F_5)\cong A_5$ (известный факт; транзитивных подгрупп $S_6$ с точностью до
сопряжённости всего $16$). Группы $A_6$ и $\mathrm{PSL}_2(\F_5)\subset A_6$ исключены, так как $G_f\not\subset A_6$. Остаются $S_6$ и
$\mathrm{PGL}_2(\F_5)$. Цикловые типы элементов $\mathrm{PGL}_2(\F_5)$ на $\F_5\cup\{\infty\}$ (с числом элементов):
$1^6$ (1), $1^22^2$ (15), $1^24$ (30), $1\,5$ (24), $2^3$ (10), $3^2$ (20), $6$ (20). Ровно эти шесть нетривиальных типов
и встретились в шаге 1; транспозиций, тройных циклов, типов $1\,2\,3$, $2\,4$ и т.\,д. в $\mathrm{PGL}_2(\F_5)$ нет. Если бы
$G_f=S_6$, то по теореме плотности Чеботарёва для доли $\frac{15}{720}$ простых $f$ по модулю $p$ раскладывался бы на
множители степеней $1,1,1,1,2$; то, что для многих сотен простых такого не происходит, делает ответ
$\mathrm{PGL}_2(\F_5)$ правдоподобным, но не доказывает его.

\emph{Шаг 3: резольвента.} Доказательство даёт резольвента для подгруппы $\mathrm{PGL}_2(\F_5)\subset S_6$. Назовём
\emph{синтемой} разбиение шести корней $\alpha_1,\dots,\alpha_6$ на три пары (синтем $15$), а \emph{тоталом}~--- набор из пяти синтем,
в которых встречаются все $15$ пар. Тоталов шесть, $S_6$ переставляет их транзитивно, и стабилизатор тотала~--- транзитивная
подгруппа $S_6$ порядка $120$, сопряжённая $\mathrm{PGL}_2(\F_5)$ (это проверяется прямым перебором). Абстрактно
стабилизатор $P_T$ любого тотала $T$ изоморфен $S_5$. Действительно, $|P_T|=720/6=120$ (формула орбит). Действие $S_6$ на шести
тоталах точно: его ядро~--- нормальная подгруппа $S_6$, т.\,е.\ $1$, $A_6$ или $S_6$, а факторгруппа по ядру действует на
шести тоталах транзитивно и потому имеет порядок $\ge6$; значит, ядро тривиально. Поэтому $P_T$ точно действует на пяти
остальных тоталах (элемент $P_T$, неподвижный на них, неподвижен на всех шести и равен $1$), и получаем вложение
$P_T\hookrightarrow S_5$, которое при $|P_T|=120$ является изоморфизмом. Для тотала $T$ положим
\[
F_T=\sum_{s\in T}\ \prod_{\{j,k\}\in s}(\alpha_j-\alpha_k)^2,\qquad R(Y)=\prod_{T}\bigl(Y-F_T\bigr)
\]
(произведение по шести тоталам). Коэффициенты $R$~--- симметрические многочлены с целыми коэффициентами от корней
унитарного многочлена $f\in\Z[X]$, поэтому они целые (\rf{5.36} для кольца $\Z$). Вычисление с точностью $100$ знаков (корни $f$ находятся
численно, коэффициенты $R$ вычисляются по определению и округляются до ближайших целых; отклонение от целых~--- порядка
$10^{-84}$, т.\,е.\ заведомо меньше $\tfrac12$, так что округление однозначно) даёт
\[
R(Y)=(Y+336)\,R_1(Y),
\]
где $R_1$~--- неприводимый над $\Q$ многочлен степени $5$, и $R$ не имеет кратных корней (вычислено на компьютере).

\emph{Шаг 4: вывод.} Элемент $\sigma\in G_f$ переводит $F_T$ в $F_{\sigma T}$. Пусть $F_{T_0}=-336$. Так как $\sigma(-336)=-336$ и
числа $F_T$ попарно различны, $\sigma T_0=T_0$; значит, $G_f$ содержится в стабилизаторе $P=P_{T_0}$ тотала $T_0$, $|P|=120$. По
шагу 1 $|G_f|$ делится на $60$ и делит $120$, т.\,е.\ $|G_f|\in\{60,120\}$. Пусть $|G_f|=60$. Подгруппа индекса $2$ содержит
квадраты всех элементов группы, а в $S_5$ подгруппа, порождённая квадратами, содержит все тройные циклы
($\rho=(\rho^2)^2$) и потому равна $A_5$. Значит, при изоморфизме $P\cong S_5$ (шаг 3) подгруппа $G_f$ индекса $2$ переходит в
$A_5$ и, следовательно, порождена квадратами элементов $P$. Но квадрат любой перестановки чётен, т.\,е.\ $G_f\subset A_6$,~---
противоречие. Итак, $G_f=P\cong\mathrm{PGL}_2(\F_5)\cong S_5$. Неприводимость $R_1$ согласуется с этим: $G_f$ транзитивно
переставляет пять остальных тоталов, и это действие задаёт изоморфизм $G_f\to S_5$ (шаг 3). \qed

% ---------------------------------------------------------------- 4-7
\begin{exercise}[4-7]
Пусть $f=X^5+aX+b$, $a,b\in\Q$. Покажите, что $G_f\cong D_5$ \textup(группа диэдра порядка $10$\textup) тогда и только тогда, когда
\begin{enumerate}[label=(\alph*)]
\item $f$ неприводим в $\Q[X]$,
\item дискриминант $D(f)=4^4a^5+5^5b^4$ многочлена $f$~--- квадрат в $\Q$,
\item уравнение $f(X)=0$ разрешимо в радикалах.
\end{enumerate}
\end{exercise}

\noindent\emph{Решение.} Формула для дискриминанта~--- случай $n=5$ формулы \rf{4.39} для $X^n+aX+b$. Будем пользоваться
описанием групп с нормальной подгруппой простого порядка из \rf{3.30a}: \emph{если транзитивная группа $G$ перестановок
пятиэлементного множества содержит нормальную подгруппу порядка $5$, то при подходящей нумерации множества элементами
$\F_5$ группа $G$ состоит из отображений $x\mapsto ux+v$ и содержит все сдвиги $x\mapsto x+v$} (\rf{3.30a}, (iii)$\Rightarrow$(iv)).
Найдём чётность таких отображений. Сдвиг $x\mapsto x+1$~--- цикл длины $5$, он чётен; $x\mapsto-x$~--- это $(1\,4)(2\,3)$,
чётно; $x\mapsto2x$~--- это $(1\,2\,4\,3)$, нечётно. Так как $2$ порождает $\F_5^\times$, а знак~--- гомоморфизм,
\[
\{x\mapsto ux+v\}\cap A_5=\{x\mapsto\pm x+v\}=:D,
\]
и $D$~--- группа порядка $10$, изоморфная $D_5$ (сдвиги~--- повороты правильного пятиугольника с вершинами $\F_5$, а
$x\mapsto-x+v$~--- отражения).

$\Rightarrow$. Пусть $G_f\cong D_5$. (a) Если $f$ приводим, то все его различные неприводимые множители $f_1,\dots,f_k$ имеют
степень $\le4$; они сепарабельны (так как $\operatorname{char}\Q=0$) и попарно взаимно просты, а поле разложения $f$
совпадает с полем разложения $f_1\cdots f_k$ (корни те же). По \rf{3.27e} $G_f$ вкладывается в $G_{f_1}\times\dots\times G_{f_k}$,
порядок которой делит $\prod_i(\deg f_i)!$ и потому не делится на $5$. А $5\mid10=|G_f|$. Значит, $f$ неприводим (и
сепарабелен). (b) В $G_f$ есть элемент порядка $5$ (Коши), т.\,е.\ цикл длины $5$, и порождённая им подгруппа индекса $2$
нормальна. По сказанному выше $G_f$ состоит из отображений $x\mapsto ux+v$ и содержит все пять сдвигов. Сопоставление
$(x\mapsto ux+v)\mapsto u$~--- гомоморфизм в циклическую группу $\F_5^\times$ с ядром из пяти сдвигов, поэтому при $|G_f|=10$
его образ~--- подгруппа порядка $2$, т.\,е.\ $\{\pm1\}$: $G_f=D$. Все её элементы чётны, $G_f\subset A_5$, и по \rf{4.2} $D(f)$~--- квадрат. (c) Группа $D_5$ разрешима
(нормальная подгруппа поворотов циклическая, факторгруппа порядка $2$), поэтому по теореме Галуа \rf{5.34} $f$ разрешим в
радикалах.

$\Leftarrow$. Пусть выполнены (a)--(c). По (a) $G_f$ транзитивна (\rf{3.27d}) и $5\mid|G_f|$. По (c) и \rf{5.34} $G_f$ разрешима, и по
\rf{3.30a} ((v)$\Rightarrow$(iv)) при подходящей нумерации корней $G_f\subset\{x\mapsto ux+v\}$ и содержит все сдвиги. По (b) и
\rf{4.2} $G_f\subset A_5$, т.\,е.\ $G_f\subset D$. Подгруппы $D$, содержащие пять сдвигов,~--- это группа сдвигов $C_5$ и сама $D$.
Остаётся исключить $G_f=C_5$. Производная $f'=5X^4+a$ имеет не больше двух вещественных корней, поэтому по теореме Ролля у $f$
не больше трёх вещественных корней (все корни $f$ простые). Значит, у $f$ есть невещественный корень, комплексное
сопряжение задаёт нетривиальный элемент порядка $2$ группы $G_f$ (поле разложения $f$ в $\CC$ переходит в себя), и
$|G_f|$ чётен. Итак, $G_f=D\cong D_5$. \qed

\begin{remark}
(1) Попутно доказано, что для трёхчленов $X^5+aX+b$ группа $C_5$ невозможна, а $D_5$ встречается: например, для
$X^5-5X+12$ (\rf{2.4d}, \rf{4.18g}(1), $D=2^{12}5^6$) и для $X^5+20X+32$, $X^5+11X+44$ (вычислено в PARI).

(2) Ни одно из условий нельзя отбросить (\rf{4.18g}): для $X^5+20X+16$ выполнены (a) и (b), но $G_f=A_5$; для $X^5+15X+12$
выполнены (a) и (c), но $G_f=F_{20}$ и $D(f)=2^{10}3^45^5$ не квадрат; для $X^5-X-1$ (\rf{4.30}) выполнено лишь (a).

(3) На практике условие (c) проверяют резольвентой Кэли (\rf{4.18e}; для трёхчленов~--- явная формула \rf{4.18f}).
\end{remark}

% ---------------------------------------------------------------- 4-8
\begin{exercise}[4-8]
Пусть $n=\prod_{i=1}^kp_i^{r_i}$, где $p_i$~--- различные простые числа. Покажите, что многочлен $f\in\F_p[X]$ степени $n$
неприводим над $\F_p$ тогда и только тогда, когда \textup{(a)} $\gcd\bigl(f(X),X^{p^{n/p_i}}-X\bigr)=1$ для всех $1\le i\le k$
и \textup{(b)} $f$ делит $X^{p^n}-X$ \textup(тест неприводимости Рабина\textup).
\end{exercise}

\noindent\emph{Указания.} Основной факт: по \rf{4.22b} многочлен $X^{p^m}-X$ равен произведению всех унитарных
неприводимых многочленов над $\F_p$, степень которых делит $m$, причём каждый входит один раз; можно считать $f$
унитарным.
\begin{enumerate}[leftmargin=*]
\item \emph{Необходимость.} Условие (b) прямо следует из основного факта. Для (a): каким может быть
$\gcd(f,X^{p^{m}}-X)$, если $f$ неприводим? Что означал бы второй вариант для степени $n$ при $m=n/p_i$?
\item \emph{Достаточность.} Из (b) выведите, что $f$ не имеет кратных неприводимых множителей и что степень $d$ каждого
неприводимого множителя $g$ делит $n$. Предположите, что $d<n$, и найдите простое $p_i$, для которого $d\mid n/p_i$
(подумайте, какое простое число делит $n/d>1$). Что тогда можно сказать о $\gcd(f,X^{p^{n/p_i}}-X)$?
\item \emph{Оба условия нужны.} Без (a): $f=X^2-X$ делит $X^{p^2}-X$. Без (b): при $n=8$ произведение неприводимых
многочленов степеней $3$ и $5$ удовлетворяет (a), так как ни $3$, ни $5$ не делят $n/2=4$. (А при $n=6$ условие (a)
уже влечёт неприводимость~--- проверьте, перебрав разбиения числа $6$.)
\item \emph{Почему тест быстрый.} Многочлены $X^{p^m}\bmod f$ вычисляются последовательным возведением в степень $p$ по
модулю $f$ ($m$ шагов, каждый~--- быстрым возведением в степень), после чего $\gcd$ находится алгоритмом Евклида;
многочлен $X^{p^n}-X$ степени $p^n$ при этом никогда не выписывается.
\end{enumerate}

% ---------------------------------------------------------------- 4-9
\begin{exercise}[4-9]
Пусть $f\in\Q[X]$~--- неприводимый многочлен, у которого есть и вещественные, и невещественные корни. Покажите, что его
группа Галуа неабелева. Можно ли отбросить условие неприводимости?
\end{exercise}

\noindent\emph{Указания.}
\begin{enumerate}[leftmargin=*]
\item Считайте поле разложения $E$ лежащим в $\CC$. Комплексное сопряжение переводит $E$ в себя (почему?) и задаёт элемент
$\sigma\in G_f$. Что $\sigma$ делает с вещественным корнем $a_3$ и с невещественным корнем $a_1$?
\item Неприводимость даёт транзитивность (\rf{3.27d}): найдётся $\tau\in G_f$ с $\tau(a_3)=a_1$. Сравните $\sigma\tau(a_3)$ и
$\tau\sigma(a_3)$.
\item Другой путь: вспомните \rf{4.9c}(b)~--- каков стабилизатор точки в транзитивной абелевой группе перестановок? Есть ли
неподвижная точка у $\sigma$?
\item Для второго вопроса рассмотрите произведение неприводимого квадратного трёхчлена с невещественными корнями на
линейный многочлен. Какова группа Галуа?
\end{enumerate}

% ---------------------------------------------------------------- 4-10
\begin{exercise}[4-10]
Пусть $F$~--- расширение Галуа поля $\Q$, и $\alpha\in F$~--- элемент, для которого класс $\alpha F^{\times2}$ не неподвижен
относительно действия $\Gal(F/\Q)$ на $F^\times/F^{\times2}$. Пусть $\alpha=\alpha_1,\dots,\alpha_n$~--- орбита $\alpha$ под
действием $\Gal(F/\Q)$. Покажите, что
\begin{enumerate}[label=(\alph*)]
\item $F[\sqrt{\alpha_1},\dots,\sqrt{\alpha_n}]/F$~--- расширение Галуа с коммутативной группой Галуа, содержащейся в
$(\Z/2\Z)^n$;
\item $F[\sqrt{\alpha_1},\dots,\sqrt{\alpha_n}]/\Q$~--- расширение Галуа с некоммутативной группой Галуа, содержащейся в
$(\Z/2\Z)^n\rtimes\Gal(F/\Q)$.
\end{enumerate}
\end{exercise}

\noindent\emph{Указания.} Положим $L=F[\sqrt{\alpha_1},\dots,\sqrt{\alpha_n}]$, $\Gamma=\Gal(F/\Q)$, и зафиксируем
квадратные корни $\beta_j$, $\beta_j^2=\alpha_j$.
\begin{enumerate}[leftmargin=*]
\item (a) Поле $L$~--- композит полей $F[\beta_j]$, каждое из которых~--- поле разложения $X^2-\alpha_j$ над $F$, т.\,е.\ расширение
Галуа степени $1$ или $2$. Воспользуйтесь \rf{3.27e} (группа Галуа произведения многочленов вкладывается в произведение
групп Галуа множителей) для $\prod_j(X^2-\alpha_j)$ или прямо: $\rho\in\Gal(L/F)$ определяется знаками $\varepsilon_j(\rho)=\pm1$
в $\rho(\beta_j)=\varepsilon_j(\rho)\beta_j$.
\item (b), \emph{Галуа.} Многочлен $\prod_j(X^2-\alpha_j)$ имеет коэффициенты в $\Q$: они неподвижны относительно $\Gamma$, так
как $\Gamma$ переставляет $\alpha_j$. Если $F$~--- поле разложения $h\in\Q[X]$, то $L$~--- поле разложения
$h\cdot\prod_j(X^2-\alpha_j)$ над $\Q$.
\item (b), \emph{вложение.} Для $\sigma\in\Gal(L/\Q)$ имеем $\sigma(\beta_j)^2=\sigma(\alpha_j)=\alpha_{s(j)}$, где $s$~---
перестановка, задаваемая $\sigma|_F$; значит, $\sigma(\beta_j)=\varepsilon_j(\sigma)\beta_{s(j)}$. Проверьте, что
$\sigma\mapsto(\varepsilon(\sigma),\sigma|_F)$~--- инъективный гомоморфизм в группу пар $(\varepsilon,\gamma)$,
$\varepsilon\in\{\pm1\}^n$, $\gamma\in\Gamma$, с умножением, в котором $\gamma$ действует на $\{\pm1\}^n$ перестановкой
координат по своему действию на орбите (<<знаковые перестановочные матрицы>>). Почему это полупрямое произведение?
\item (b), \emph{некоммутативность.} Возьмите $s\in\Gamma$ с $s(\alpha)F^{\times2}\ne\alpha F^{\times2}$. Объясните, почему тогда
$\alpha\notin F^{\times2}$ и $s(\alpha)\notin F^{\times2}$, и, используя лемму из решения упражнения 2-1 (когда $a\in F^\times$
становится квадратом в $F[\sqrt b]$), покажите, что $[F[\sqrt\alpha,\sqrt{s(\alpha)}]:F]=4$. Выберите
$\rho\in\Gal(L/F)$ с $\rho(\sqrt\alpha)=-\sqrt\alpha$, $\rho(\sqrt{s(\alpha)})=\sqrt{s(\alpha)}$ и продолжение $\sigma\in\Gal(L/\Q)$
элемента $s$. Сравните $\sigma\rho(\sqrt\alpha)$ и $\rho\sigma(\sqrt\alpha)$.
\item \emph{Пример для проверки.} $F=\Q[\sqrt2]$, $\alpha=1+\sqrt2$: орбита $\{1\pm\sqrt2\}$, отношение
$\frac{1-\sqrt2}{1+\sqrt2}=-(\sqrt2-1)^2$ не квадрат в $F\subset\R$; $L$~--- поле разложения $X^4-2X^2-1$ (ср.\ \rf{3.25c}), и его
группа $D_4$ совпадает с $(\Z/2\Z)^2\rtimes\Z/2\Z$. Обратите внимание, что в (b) утверждается лишь вложение $\Gal(L/\Q)$ в полупрямое произведение как подгруппы, а не
равенство.
\end{enumerate}
